\section{Integración compleja}\label{sec:integracion-compleja}
\subsection{Introducción a la integración compleja}

\begin{definición}[Integral de una función compleja de variable real]
    Sea $\varphi: [a, b] \to \mathbb{C}$ una función, definimos:
    \[
    \int_a^b \varphi(t) \, dt = \int_a^b \Re (\varphi(t)) \, dt + i \int_a^b \Im (\varphi(t)) \, dt
    \]
\end{definición}

\begin{proposición}[Propiedades de la integral de funciones complejas]
    Sean $\varphi, \psi: [a, b] \to \mathbb{C}$ funciones integrables y $\alpha, \beta \in \mathbb{C}$. Entonces:
    \begin{enumerate}
        \item Linealidad:
        \[
        \int_a^b (\alpha\varphi(t) + \beta\psi(t)) \, dt = 
        \alpha\int_a^b \varphi(t) \, dt + \beta\int_a^b \psi(t) \, dt
        \]
        \item Aditividad en el intervalo: Si $a < c < b$, entonces:
        \[
        \int_a^b \varphi(t) \, dt = \int_a^c \varphi(t) \, dt + \int_c^b \varphi(t) \, dt
        \]
        \item Acotación:
        \[
        \left| \int_a^b \varphi(t) \, dt \right| \leq 
        \int_a^b |\varphi(t)| \, dt
        \]
    \end{enumerate}
\end{proposición}

\begin{proof}
    Observemos que si $z \in \mathbb{C} \setminus \{0\}$ y $\theta = \arg(z)$ entonces $|z| = z e^{-i\theta}$. Por tanto, si $\varphi: [a, b] \to \mathbb{C}$ es integrable, y tenemos que $\theta_0$ es un argumento de $\int_a^b \varphi(t) \, dt$, entonces:
    \[
    \left| \int_a^b \varphi(t) \, dt \right| =
    e^{-i\theta_0} \int_a^b \varphi(t) \, dt =
    \int_a^b \varphi(t) e^{-i\theta_0} \, dt =
    \int_a^b \Re(\varphi(t) e^{-i\theta_0}) \, dt + i \int_a^b \Im \varphi(t) e^{-i\theta_0} \, dt
    \]
    Esta última integral es real, por lo que $\Im \int_a^b \varphi(t) e^{-i\theta_0} \, dt = 0$. Por tanto:
    \[
    \left| \int_a^b \varphi(t) \, dt \right| =
    \int_a^b \Re(\varphi(t) e^{-i\theta_0}) \, dt 
    \underbrace{\leq}_{\Re (z) \leq |\Re (z)| \leq |z|}
    \int_a^b |\varphi(t) e^{-i\theta_0}| \, dt =
    \int_a^b |\varphi(t)| \, dt
    \]
\end{proof}

\begin{observación}[Cambio de variable en integrales complejas]
    Si $\varphi: [a, b] \to \mathbb{C}$ es una función derivable con derivada continua y $f: \varphi([a, b]) \to \mathbb{C}$ es continua, entonces, de forma análoga al caso real, se cumple que:
    \[
    \int_{\varphi (a)}^{\varphi (b)} f(z) \, dz = \int_a^b f(\varphi(t)) \varphi'(t) \, dt
    \]
\end{observación}

\subsection{Curvas en $\mathbb{C}$}

Recordando las curvas en $\mathbb{C}$, tenemos que una curva $\gamma$ es una función continua $\gamma: [a, b] \to \mathbb{C}$.

\ejemplo{
    Veamos algunos ejemplos de curvas en $\mathbb{C}$:
    \begin{enumerate}
        \item La circunferencia unidad centrada en el origen viene dada por $\gamma: [0, 2\pi] \to \mathbb{C}$, $\gamma(t) = \cos t + i \sin t$. 
        \item La circunferencia de centro $z_0$ y radio $r$ viene dada por $\gamma: [0, 2\pi] \to \mathbb{C}$, $\gamma(t) = (z_0 + r\cos t, z_0 + r\sin t)$.
        \item Un segmento en $\mathbb{C}$ definido por los puntos $z_1, z_2 \in \mathbb{C}$ viene dado por $\gamma: [0, 1] \to \mathbb{C}$, $\gamma(t) = (1-t)z_1 + tz_2$.
    \end{enumerate}
}

\begin{definición}[Suma de curvas en $\mathbb{C}$]
    Si $\gamma, \tilde{\gamma}$ son dos curvas definidas por $\gamma: [a, b] \to \mathbb{C}$ y $\tilde{\gamma}: [c, d] \to \mathbb{C}$, y que cumplen que $\gamma(b) = \tilde{\gamma}(c)$, podemos definir la curva $\gamma \cup \tilde{\gamma} = \gamma + \tilde{\gamma}$ como:

    \[
    (\gamma + \tilde{\gamma})(t) =
    \begin{cases}
    \gamma(t) & \text{si } t \in [a, b] \\
    \tilde{\gamma}(t + (c - b)) & \text{si } t \in [b, b + (d - c)]
    \end{cases}
    \]
\end{definición}

\begin{definición}[Curva recorrida en sentido contrario]
    Si $\gamma: [a, b] \to \mathbb{C}$ es una curva en $\mathbb{C}$, podemos denotar $\gamma^-$ a la curva recorrida en sentido contrario, es decir, $\gamma^-: [a, b] \to \mathbb{C}$, $\gamma^-(t) = \gamma(a + b - t)$.
\end{definición}

\begin{definición}[Longitud de una curva en $\mathbb{C}$]
    Si $\gamma: [a, b] \to \mathbb{C}$ es una curva de clase $\mathcal{C}^1$ en $[a, b]$, podemos definir la longitud de $\gamma$ como:
    \[
    \Long(\gamma) = \int_a^b |\gamma'(t)| \, dt
    \]
    Teniendo en cuenta que $\gamma = \gamma_1 + i\gamma_2$, y pr tanto $\gamma' = \gamma_1' + i\gamma_2'$.
\end{definición}

\begin{definición}[Traza de una curva en $\mathbb{C}$]
    Si $\gamma: [a, b] \to \mathbb{C}$ es una curva en $\mathbb{C}$, definimos la traza de $\gamma$ y denotaremos por $\gamma^*$ a $\gamma([a, b])$.
\end{definición}

Si $\gamma$ es inyectiva, en $[a, b)$ a veces se identifica con su traza $\gamma^*$ y una orientación. En este caso además, la longitud de $\gamma$ coincide con la longitud de su traza $\Long(\gamma^*) = \Long(\gamma) = \int_a^b |\gamma'(t)| \, dt$.

\begin{definición}[Integral de una función compleja a lo largo de una curva]
    Sea $\gamma: [a, b] \to \mathbb{C}$ una curva de clase $\mathcal{C}^1$ en $[a, b]$ y $f: \gamma^* \to \mathbb{C}$ una función continua. Definimos la integral de $f$ a lo largo de la curva $\gamma$ como:
    \[
    \int_\gamma f(z) \, dz = \int_a^b f(\gamma(t)) \gamma'(t) \, dt
    \]
\end{definición}

\ejemplo{
    Calculemos la integral $\int_{S(0;1)} \frac{1}{z} \, dz$, donde $S(0;1) = C(0;1) = \{ z \in \mathbb{C} : |z| = 1 \}$ es la circunferencia unidad centrada en el origen. Para ello, tomamos la parametrización $\gamma: [0, 2\pi] \to \mathbb{C}$, $\gamma(t) = e^{it} = \cos t + i \sin t$. Entonces:
    \[
    \int_{S(0;1)} \frac{1}{z} \, dz = 
    \int_0^{2\pi} \frac{1}{\gamma(t)} \gamma'(t) \, dt =
    \int_0^{2\pi} \frac{-\sin t + i \cos t}{\cos t + i \sin t} \, dt =
    \]
    Multiplicando por el conjugado del denominador:
    \[
    = \int_0^{2\pi} \frac{(-\sin t + i \cos t) \overline{(\cos t + i \sin t)}}{|\cos t + i \sin t|^2} \, dt =
    \int_0^{2\pi} (-\sin t + i \cos t)(\cos t - i \sin t) \, dt =
    \]
    \[
    = \int_0^{2\pi} -\sin t \cos t + \cos t \sin t + i(\cos^2 t + \sin^2 t) \, dt =
    \int_0^{2\pi} i \, dt = 2\pi i
    \]
    \[
        \gamma : [0, 2\pi] \to \mathbb{C}, \quad \gamma(t) = e^{it} = \cos t + i \sin t
    \]
    \[
        \gamma'(t) = ie^{it} = -\sin t + i \cos t
    \]
    Luego
    \[
        \int_{S(0;1)} \frac{1}{z} \, dz = 
        \int_0^{2\pi} \frac{1}{\gamma(t)} \gamma'(t) \, dt =
        \int_0^{2\pi} \frac{1}{e^{it}} ie^{it} \, dt =
    \]
    \[
        = \int_0^{2\pi} i \, dt = 2\pi i
    \]
}

\ejemplo{

    \[
        \int_{[2i,2]} z^2 \, dz
    \]
    \[
        \gamma : [0, 1] \to \mathbb{C}, \quad \gamma(t) = (1-t)2i + t2 = 2t + 2i(1-t)
    \]
    \[
        \gamma'(t) = 2 - 2i
    \]
    \[
        \int_{[2i,2]} z^2 \, dz = 
        \int_0^1 (\gamma(t))^2 \gamma'(t) \, dt =
        \int_0^1 (2t + 2i(1-t))^2 (2 - 2i) \, dt =
    \]
    \[
        = (2-2i) \int_{0}^1 [(2-2i)^2t^2-4+4i(2-2i)t] \, dt =
        (2-2i) \left[ (2-2i)^2 \frac{1}{0}t^2-4\int_{0}^{1} dt + 4i(2-2i) \frac{1}{2}t \right] = ... = \frac{8}{3}(1+i)
    \]

}

\subsection{Propiedades de la integral compleja}

\begin{proposición} [Propiedades de la integral]
    Sea $\gamma: [a, b] \to \mathbb{C}$ una curva de clase $\mathcal{C}^1$ en $[a, b]$ y $f, g: \gamma^* \to \mathbb{C}$ funciones continuas. Entonces:
    \begin{enumerate}
        \item Linealidad:
        \[  
        \int_\gamma (\alpha f(z) + \beta g(z)) \, dz =
        \alpha \int_\gamma f(z) \, dz + \beta \int_\gamma g(z) \, dz
        \]
        \item Si $c \in (a, b)$, y llamamos $\gamma_1 = \gamma|_{[a, c]}$ y $\gamma_2 = \gamma|_{[c, b]}$, entonces:
        \[
        \int_\gamma f(z) \, dz = \int_{\gamma_1} f(z) \, dz + \int_{\gamma_2} f(z) \, dz
        \]
        \item \[
            \int_{\gamma^-} f(z) \, dz = - \int_\gamma f(z) \, dz
        \]
        \item Si $\varphi: [c, d] \to [a, b]$ es una función biyectiva y de clase $\mathcal{C}^1$ en $[c, d]$ con $0 < \varphi'$ en $[c, d]$, y definimos $\tilde{\gamma} = \gamma \circ \varphi: [c, d] \to \mathbb{C}$, entonces:
        \[
        \int_{\tilde{\gamma}} f(z) \, dz = \int_\gamma f(z) \, dz
        \]
    \end{enumerate}
\end{proposición}

\begin{proof}
    \leavevmode
    \begin{enumerate}
        \item Inmediata.
        \item \[
        \int_\gamma f(z) \, dz = \int_a^b f(\gamma(t)) \gamma'(t) \, dt = \int_{a}^c f(\gamma(t)) \gamma'(t) \, dt + \int_c^b f(\gamma(t)) \gamma'(t) \, dt = \int_{\gamma_1} f(z) \, dz + \int_{\gamma_2} f(z) \, dz
        \]
        \item Tenemos que $\gamma^-: [a, b] \to \mathbb{C}$, manda $t \mapsto \gamma(a + b - t)$. Por tanto:
        \[
            \gamma'^-(t) = -\gamma'(a + b - t)
        \]
        Así,
        \[
            \int_{\gamma^-} f(z) \, dz = \int_a^b f(\gamma^-(t)) \gamma'^-(t) \, dt = \int_a^b f(\gamma(a + b - t)) (-\gamma'(a + b - t)) \, dt
        \]
        \[
            \underbrace{=}_{s = a + b - t \quad ds = -dt} \int_b^a f(\gamma(s)) (-1) \gamma'(s) (-ds) = -\int_a^b f(\gamma(s)) \gamma'(s) \, ds = -\int_\gamma f(z) \, dz
        \]
        \item \[
            \int_{\tilde{\gamma}} f(z) \, dz = \int_c^d f(\tilde{\gamma}(t)) \tilde{\gamma}'(t) \, dt = \int_c^d f(\gamma(\varphi(t))) \gamma'(\varphi(t)) \varphi'(t) \, dt
        \]
        \[
            \underbrace{=}_{s = \varphi(t) \quad ds = \varphi'(t) dt} \int_{a}^b f(\gamma(s)) \gamma'(s) \, ds = \int_\gamma f(z) \, dz
        \]
    \end{enumerate}
\end{proof}

\begin{observación}
    (3) nos dice que $\int_\gamma f(z) \, dz$ es una integral orientada (si cambiamos la orientación de la curva, cambiamos el signo de la integral).\\
    (4) nos dice que $\int_\gamma f(z) \, dz$ no depende de la parametrización de la curva, (esto es, si damos otra parametrización a la misma curva, el valor de la integral no cambia).\\
    Ojo: no se cumple $\left| \int_{\gamma} f(z) \, dz \right| \leq \int_{\gamma} |f(z)| \, dz$, ya que la integral de la derecha es un número complejo.\\

\end{observación}    

Pero sí tenemos la siguiente acotación en las hipotésis anteriores
    \[
        \left| \int_{\gamma} f(z) \, dz \right| \leq M \, \Long(\gamma^*) \quad \text{si } |f(z)| \leq M \text{ para todo } z \in \gamma^*
    \]

En efecto,

\[
    \left| \int_{\gamma} f(z) \, dz \right| = \left| \int_a^b f(\gamma(t)) \gamma'(t) \, dt \right| \leq \int_a^b |f(\gamma(t))| |\gamma'(t)| \, dt \leq \int_a^b M|\gamma'(t)| \, dt = M \int_{a}^b |\gamma'(t)| \, dt = M \, \Long(\gamma^*)
\]

\ejemplo{

    Acota $\int_\gamma \frac{1}{z^4+1} \, dz$ donde $\gamma$ es la circunferencia de centro $0$ y radio $R > 1$.
    Recordamos que si no damos la parametrización de la curva, se entiende que es la parametrización positiva (sentido antihorario).

    \[
        \gamma : [0, 2\pi] \to \mathbb{C}, \quad \gamma(t) = Re^{it} = R(\cos t + i \sin t)
    \]
    Por tanto,
    \[  
        \gamma'(t) = iRe^{it} = R(-\sin t + i \cos t)
    \]
    Así,
    \[
        \left| \int_\gamma \frac{1}{z^4+1} \, dz \right| = \left| \int_0^{2\pi} \frac{1}{\gamma(t)^4 + 1} \gamma'(t) \, dt \right| = \left| \int_0^{2\pi} \frac{1}{R^4 e^{4it} + 1} iRe^{it} \, dt \right| 
    \]
    \[
        \leq \int_0^{2\pi} \left| \frac{1}{R^4 e^{4it} + 1} iRe^{it} \right| \, dt = \int_0^{2\pi} \frac{R}{|R^4 e^{4it} + 1|} \, dt 
    \]
    \[
        \underbrace{\leq}_{|R^4 e^{4it} + 1| \geq |R^4e^{4it}| - |1| = R^4 - 1} \int_0^{2\pi} \frac{R}{R^4 - 1} \, dt = \frac{2\pi R}{R^4 - 1}
    \]

}

Si $\gamma : [a, b] \to \mathbb{C}$ es una curva de clase $\mathcal{C}^1$ a trozos (es decir, existe una partición finita $\{t_0, t_1, \ldots, t_n\}$ de $[a, b]$ donde $a = t_0 < t_1 < \ldots < t_n = b$ tal que $\gamma_i = \gamma|_{[t_{i-1}, t_i]}$ es de clase $\mathcal{C}^1$ en $[t_{i-1}, t_i]$ para todo $i = 1, \ldots, n$), definimos la integral 

\[
    \int_\gamma f(z) \, dz = \sum_{i=1}^n \int_{\gamma_i} f(z) \, dz
\]

Se puede comprobar que esta definición no depende de la partición escogida.

\subsection{Teorema de Barrow y existencia de primitivas}

Para esta integral tenemos Regla de Barrow.

\begin{teorema} [Regla de Barrow]
    Sea $\gamma: [a, b] \to \mathbb{C}$ una curva de clase $\mathcal{C}^1$ en $[a, b]$ y $f: \gamma^* \to \mathbb{C}$ una función continua. Si existe $F: \gamma^* \to \mathbb{C}$ derivable en $\gamma^*$ tal que $F'(z) = f(z)$ para todo $z \in \gamma^*$, entonces:
    \[
        \int_\gamma f(z) \, dz = F(\gamma(b)) - F(\gamma(a))
    \]
\end{teorema}

\textbf{Nota:} Que $F$ sea derivable en $\gamma^*$ significa que $F$ es derivable en un abierto que contiene a $\gamma^*$.

\begin{proof}
    \[
        \int_\gamma f(z) \, dz = \int_a^b f(\gamma(t)) \gamma'(t) \, dt = \int_a^b F'(\gamma(t)) \gamma'(t) \, dt = \int_a^b (F \circ \gamma)'(t) \, dt
    \]
    \[
    = \int_a^b \Re ((F \circ \gamma)'(t)) \, dt + 
    i \int_a^b \Im ((F \circ \gamma)'(t)) \, dt =^*
    \int_a^b (\Re (F \circ \gamma))'(t) \, dt +
    i \int_a^b (\Im (F \circ \gamma))'(t) \, dt \underbrace{=}_{\text{Regla de Barrow en } \mathbb{R}}
    \]
    \[
    = \left[ \Re (F \circ \gamma)(t) \right]_a^b + i \left[ \Im (F \circ \gamma)(t) \right]_a^b =
    \Re (F \circ \gamma(b)) - \Re (F \circ \gamma(a)) + i(\Im (F \circ \gamma(b)) - \Im (F \circ \gamma(a))) =
    \]
    \[
    =F \circ \gamma(b) - F \circ \gamma(a) = F(\gamma(b)) - F(\gamma(a))
    \]
    Usando en $*$ que si $\phi: [a, b] \to \mathbb{C}$ es derivable y se define como $\phi = \phi_1 + i\phi_2$, entonces $\phi' = \phi_1' + i\phi_2'$ y por tanto $\Re (\phi'(t)) = (\Re \phi (t))'$ y $\Im (\phi'(t)) = (\Im \phi (t))'$.
\end{proof}

\ejemplo{
    \[
        \int_{[2i,2]} z^2 \, dz = F(2) - F(2i) \quad \text{si } F(z) = \frac{z^3}{3} \quad \text{ya que } F'(z) = z^2
    \]
    \[
        = \frac{1}{3}(2^3 - (2i)^3) = \frac{8}{3}(1+i)
    \]
}

\begin{definición}[Curva cerrada]
    Una curva $\gamma: [a, b] \to \mathbb{C}$ es cerrada si $\gamma(a) = \gamma(b)$.
\end{definición}

\begin{corolario}[Integral en curva cerrada]
    Sea $\Omega \subset \mathbb{C}$ un abierto y $f: \Omega \to \mathbb{C}$ una función continua en $\Omega$. Si existe $F: \Omega \to \mathbb{C}$ derivable en $\Omega$ tal que $F'(z) = f(z)$ para todo $z \in \Omega$, entonces, para toda curva cerrada $\gamma$ contenida en $\Omega$ se cumple que:
    \[
        \int_\gamma f(z) \, dz = 0
    \]
\end{corolario}

De esto podemos afirmar que no existe una determinacion continua del logaritmo en $\mathbb{C} \setminus \{0\}$, ya que si existiera seria holomorfa en $\mathbb{C} \setminus \{0\}$ y su derivada seria $\frac{1}{z} \quad \forall z \in \mathbb{C} \setminus \{0\}$. Y por tanto si tomamos la integral a lo largo de cualquier curva cerrada $\mathcal{C}^1$ de su derivada deberia ser $0$, pero ya hemos visto que no es así, por ejemplo, si tomamos la circunferencia unidad centrada en el origen:
\[
\int_{C(0;1)} \frac{1}{z} \, dz = 2\pi i \neq 0
\]

Ahora nos podemos preguntar cuándo una función $f: \Omega \to \mathbb{C}$ admite una primitiva en $\Omega$.

\begin{definición}[Camino en $\mathbb{C}$]
    Sea $\Omega \subset \mathbb{C}$ un abierto. Un camino en $\Omega$ es una curva $\gamma: [a, b] \to \Omega$ de clase $\mathcal{C}^1$ ó $\mathcal{C}^1$ a trozos en $[a, b]$.
\end{definición}

\begin{definición}[Dominio o region]
    Diremos que $\Omega \subset \mathbb{C}$ es un dominio o región si $\Omega$ es un abierto y conexo.
\end{definición}

\begin{teorema}[Teorema de existencia de primitivas]
    Sea $\Omega \subset \mathbb{C}$ un dominio y $f: \Omega \to \mathbb{C}$ una función continua en $\Omega$. Entonces son equivalentes:
    \begin{enumerate}
        \item $f$ admite una primitiva en $\Omega$.
        \item $\exists F \in \mathcal{H}(\Omega)$ tal que $F' = f$ en $\Omega$.
        \item Para todo camino cerrado $\gamma$ en $\Omega$ se cumple que:
        \[
            \int_\gamma f(z) \, dz = 0
        \]
    \end{enumerate} \label{teo:existencia-primitivas}
\end{teorema}

\begin{proof}
    Claramente $(1) \Leftrightarrow (2)$ por definición de primitiva.\\
    Veamos entonces las implicaciones con $(3)$:
    \begin{enumerate}
        \item[$(1) \implies (3)$] Esto se cumple claramente por la Regla de Barrow.
        \item[$(1) \impliedby (3)$] Supongamos que $\int_{\gamma} f(z) \, dz = 0 \quad \forall \gamma \subset \Omega$. Esto implica que si $\gamma_1, \gamma_2$ son dos curvas $\mathcal{C}^1$ ó $\mathcal{C}^1$ a trozos en $\Omega$ con $\gamma_1(a) = \gamma_2(a)$ y $\gamma_1(b) = \gamma_2(b)$, entonces:
        \[
            \int_{\gamma_1} f(z) \, dz = \int_{\gamma_2} f(z) \, dz
        \]
        Ya que $\gamma_1 + \gamma_2^-$ es una curva cerrada en $\Omega$, y por tanto:
        \[
            0 = \int_{\gamma_1 + \gamma_2^-} f(z) \, dz = 
            \int_{\gamma_1} f(z) \, dz + \int_{\gamma_2^-} f(z) \, dz = \int_{\gamma_1} f(z) \, dz - \int_{\gamma_2} f(z) \, dz
        \]
        Fijemos entonces un punto $z_0 \in \Omega$ y definamos $F: \Omega \to \mathbb{C}$ como:
        \[
            \begin{array}{l}
                F: \Omega \to \mathbb{C} \\
                z \mapsto \int_{\gamma_{z_0, z}} f(w) \, dw
            \end{array}
        \]
        Siendo $\gamma_{z_0, z}$ una poligonal en $\Omega$ que une $z_0$ con $z$, que existe al ser $\Omega$ un dominio, por lo que es conexo por poligonales.\\
        Por lo que $F$ está bien definida, ya que si $\tilde{\gamma}_{z_0, z}$ es otra poligonal en $\Omega$ que une $z_0$ con $z$, entonces ambas tienen los mismos extremos, y por tanto coinciden sus integrales como hemos visto antes.\\
        Veamos que $F'=f$ en $\Omega$.\\
        Sea $z_1 \in \Omega$, tenemos que ver que $\lim_{h \to 0} \frac{F(z_1 + h) - F(z_1)}{h} = f(z_1)$.\\
        Como $z_1 \in \Omega$ y $\Omega$ es abierto $\implies \exists r > 0$ tal que $D(z_1; r) \subset \Omega$.\\
        Por tanto, si $0 < |h| < r$, tenemos que:
        \[
        F(z_1 + h) - F(z_1) = 
        \int_{\gamma_{z_0, z_1 + h}} f(w) \, dw -
        \int_{\gamma_{z_0, z_1}} f(w) \, dw =
        \int_{\gamma_{z_1, z_1 + h}} f(w) \, dw
        \]
        Luego:
        \[
        \left| \frac{F(z_1 + h) - F(z_1)}{h} - f(z_1) \right| =
        \left| \frac{F(z_1 + h) - F(z_1) - hf(z_1)}{h} \right| \underbrace{=}_{hf(z_1)=\int_{\gamma_{z_1, z_1 + h}} f(z_1) \, dz}
        \]
        \[
        = \frac{1}{|h|} \left| \int_{\gamma_{[z_1, z_1 + h]}} f(w) dw - \int_{\gamma_{[z_1, z_1 + h]}} f(z_1) \, dz \right| 
        \]
        \[
        = \frac{1}{|h|} \left| \int_{[z_1, z_1 + h]} (f(w) - f(z_1)) \, dw \right| \leq \frac{1}{|h|} \left(\max_{w \in [z_1, z_1 + h]} |f(w) - f(z_1)| \right) \Long([z_1, z_1 + h]) 
        \]
        \[
            = \max_{w \in [z_1, z_1 + h]} |f(w) - f(z_1)| \xrightarrow[h \to 0]{} 0 \text{ por la continuidad de } f \text{ en } z_1
        \]
        pues dado $\varepsilon > 0$, existe $\delta > 0$ tal que si $|\tilde{h} | < \delta \implies |f(z_1 + \tilde{h}) - f(z_1)| < \varepsilon$. Luego, $F$ es derivable en $z_1$ y $F'(z_1) = f(z_1)$.
    \end{enumerate}
\end{proof}

Ya sabemos que no toda función holomorfa en un dominio admite una primitiva en él, por ejemplo, $f(z) = \frac{1}{z}$ en $\Omega = \mathbb{C} \setminus \{0\}$, siendo $f(z)$ holomorfa en $\Omega$, pero $\int_{C(0;1)} \frac{1}{z} \, dz = 2\pi i \neq 0$.\\